\subsection{The actual image has the required long-root sign}
\label{sec:kt-actual-image-polarization}

The sign issue in Lemma~6.8 of Berger--Klosin,
arXiv:2611.19007v1, can be repaired for the matrix-image algebra
of the stipulated normalized source.  This does not identify that
algebra with the universal Cayley--Hamilton quotient in the
preprint.  We give a direct arithmetic injection for the actual
image, keeping this distinction explicit.  The remaining adjacent
rational Selmer group and the main vanishing gate are unchanged.

Write $H=G_E\subset G_{\mathbf Q}$, choose a complex conjugation
$c$, and put $g^*=cg^{-1}c$ for $g\in H$.  Matrix indices in this
section are $1,2,3$.  Set
\[
 J=\begin{pmatrix}0&0&1\\0&1&0\\1&0&0\end{pmatrix},\qquad
 \tau(X)=JX^tJ.                                               \tag{AP.1}
\]
This is a conjugate-self-duality pairing.  It does not say that
$\rho(g)$ preserves $J$ for every $g\in H$ without conjugating
the Galois argument.

\begin{lemma}[KT-ACTUAL-IMAGE-SYMMETRIC-POLARIZATION]
Let $R$ be a commutative local ring with $2\in R^\times$, and
let $r:G_{\mathbf Q}\to GG_3(R)$ be a normalized marked lift
in the source's conventions.  Put $\rho=r|_H$ and
$\mathcal E=R[\rho(H)]\subset M_3(R)$.  Then
\[
 \rho(g^*)=\tau(\rho(g)),\qquad
 \rho(g^c)^tJ\rho(g)=J.                                     \tag{AP.2}
\]
The standard diagonal idempotents belong to $\mathcal E$.
Writing
\[
 e_i\mathcal E e_j=\mathfrak a_{ij}E_{ij},\qquad
 \mathfrak a_{ii}=R,                                        \tag{AP.3}
\]
the anti-involution induces the identity on
$\mathfrak a_{13}E_{13}$ and $\mathfrak a_{31}E_{31}$, and
$\mathfrak a_{12}=\mathfrak a_{23}$,
$\mathfrak a_{21}=\mathfrak a_{32}$ as ideals of $R$.
These identities persist under coefficient specialization,
including nonreduced specializations.  At a literal Type~II
point the outer direct summand has the perfect symmetric
pairing $\left(\begin{smallmatrix}0&1\\1&0\end{smallmatrix}\right)$.
\end{lemma}

\begin{proof}
The source defines $\theta(M)=JM^{-t}J$ and
$GG_3=\mathrm{GL}_3\rtimes\{\pm1\}$.  Normalization says
$r(c)=(1,-1)$.  The semidirect-product multiplication gives
the first equality in (AP.2); inversion gives the second.
Both are identities over $R$, and $J^t=J$, $J^2=1$.

For the marked element $g_1\in H$ the source proves
\[
 z:=\rho(g_1)-\rho(g_1^*)
       =\operatorname{diag}(\zeta,0,-\zeta),\qquad
 \zeta\in R^\times.                                        \tag{AP.4}
\]
Its reduction is $\Psi(g_1)-\Psi(g_1)^{-1}$; the marking
requires $\Psi(g_1)^2-1$ to be a unit.  Thus
\[
 e_1=\frac{z^2+\zeta z}{2\zeta^2},\qquad
 e_2=1-\frac{z^2}{\zeta^2},\qquad
 e_3=\frac{z^2-\zeta z}{2\zeta^2}                            \tag{AP.5}
\]
belong to $\mathcal E$ and are precisely $E_{11},E_{22},E_{33}$.
Consequently (AP.3) is an actual matrix GMA, with ideals in
$R$, rather than a choice of fractional coordinates for a
pseudorepresentation.  Since
$\tau(E_{ij})=E_{4-j,4-i}$ and $\tau$ fixes the coefficients,
the assertions about its blocks follow immediately.

Every coefficient map carries $J$ to the same invertible
symmetric matrix and preserves (AP.2)--(AP.5).  For Type~II,
the source's definition is exactly the outer corner plus the
middle block.  The middle line has pairing $1$, its orthogonal
complement is the outer free rank-two summand, and the
restricted pairing has determinant $-1$.  Its perfectness
therefore survives specialization.  This proves the sign on
that summand directly; it is not inferred from odd ambient
dimension.
\end{proof}

\paragraph{The stable lattice is supplied by the actual source.}
For an actual specialization $f:S.R\to\overline{\mathbf Q}_p$
with \texttt{IsLocPt}, the checked finite-coefficient argument
in \texttt{TanDet.lean:182--216} makes $S.R/\ker f$ finite
over the base coefficient DVR.  Its fraction field is a
finite extension $F/\mathbf Q_p$.  Source integrality applies
to $g$ and $g^{-1}$, so the matrices lie in
$\mathrm{GL}_3(\mathcal O_F)$.  The lattice
$\mathcal O_F^3$ has the perfect symmetric pairing $J$.
No arbitrary stable lattice is replaced by a self-dual one
in this argument.  In particular, a later lattice selected
by a different extension theorem cannot silently be assumed
to retain this integral pairing.

\begin{theorem}[KT-ACTUAL-LONG-ROOT-SELMER-INJECTION]
Let $p>3$, let $E=\mathbf Q(\sqrt D)$ be imaginary quadratic
with $p\nmid D$, and put $\eta=\chi_D$.
Let $e\ge2$ be even with $(p-1)\nmid2e$.
Let $\mathcal O$ be a finite coefficient DVR with residue
field $k$ and fraction field $F$.  Suppose that
$\rho:H\to\mathrm{GL}_3(\mathcal O)$ is the restriction of
a marked normalized source lift, with actual residual matrix
\[
 \overline\rho=\operatorname{diag}(\psi,1,\psi^{-1}),
 \qquad\psi=\omega^e|_H.                                   \tag{AP.6}
\]
Assume that $\rho$ is unramified away from $p$ over $E$ and,
at every $v\mid p$, has a decomposition-group stable ordered
flag whose inertia characters are
$(\varepsilon_p^e,1,\varepsilon_p^{-e})$.
For the actual algebra
$\mathcal E=\mathcal O[\rho(H)]$ and (AP.3), there is an
injective $k$-linear map
\[
 \operatorname{Hom}_{\mathcal O}
   (\mathfrak a_{31}/\mathfrak a_{32}\mathfrak a_{21},k)
 \lhook\joinrel\longrightarrow
 \mathcal S_p\bigl(k(\eta\omega^{-2e})\bigr),               \tag{AP.7}
\]
where the right side is zero on $G_{\mathbf Q_p}$ and on
every inertia group $I_q$ for $q\ne p$.
This theorem applies to actual Type~I source specializations
for $N=2$, $\psi_{\rm finite}=1$ under GD and the carried
closed source convention.
\end{theorem}

\begin{proof}
All the ideals in (AP.3) are finite $\mathcal O$-modules:
$\mathcal E$ is a submodule of the finite module
$M_3(\mathcal O)$ and $\mathcal O$ is Noetherian.
For $f$ in the left side of (AP.7), the formula
\[
 \Pi_f(X)=
 \begin{pmatrix}
  \overline X_{33}&f(X_{31})\\0&\overline X_{11}
 \end{pmatrix}\qquad(X\in\mathcal E)                       \tag{AP.8}
\]
is an algebra homomorphism.  Indeed the reduction of
$\mathcal E$ is diagonal by (AP.6), while the middle term
$X_{32}Y_{21}$ in $(XY)_{31}$ is annihilated by $f$.
Thus (AP.8) restricted to $H$ is an extension of $\psi$
by $\psi^{-1}$, with cocycle
\[
             b_f(g)=\psi(g)^{-1}f(\rho(g)_{31}).           \tag{AP.9}
\]
This is the rank-one instance of Bella\"iche--Chenevier,
Theorem~1.5.5, whose matrix construction and injectivity
proof apply here without a choice of a new lattice.
For completeness, if this $H$-extension splits, the same
change of basis splits (AP.8) on its $\mathcal O$-linear
span $\mathcal E$.  Its restrictions to $e_1,e_3$ are
already the two standard diagonal projectors.  The splitting
therefore forces $f$ to vanish on $e_3\mathcal E e_1$.
The entries $\rho(g)_{31}$ span $\mathfrak a_{31}$, so
$f=0$.  Continuity follows because these are finite modules
over a complete DVR and $\rho$ is continuous.

The sign is exact at the cocycle level.  By (AP.2) and the
identity action on the long block,
\[
 \rho(g^c)_{31}=\rho(g^{-1})_{31},\qquad
 b_f(g^c)=-b_f(g).                                         \tag{AP.10}
\]
For the second equality, use
$b_f(g^{-1})=-\psi(g)^2b_f(g)$ and
$\psi(g^c)=\psi(g)$.  Since $e$ is even, the natural action
of $c$ on $k(\omega^{-2e})$ is trivial.  Hence this is the
negative-conjugation eigenspace.  Restriction through the
index-two subgroup $H$ identifies it with
$H^1(\mathbf Q,k(\eta\omega^{-2e}))$; the order two is
invertible in $k$.

Here is the local flag comparison, including its integral
coordinates.  At $v\mid p$ choose an inertia element $t$
whose three residual eigenvalues are distinct.  This is
possible because $p\nmid D$ and $(p-1)\nmid2e$.
The actual eigenvalues of $\rho(t)$ are
$\lambda_1=\varepsilon_p(t)^e$, $\lambda_2=1$,
$\lambda_3=\varepsilon_p(t)^{-e}$, with unit differences.
The spectral projectors
\[
 q_i=\prod_{j\ne i}
       \frac{\rho(t)-\lambda_j}{\lambda_i-\lambda_j}
       \quad\hbox{belong to }\mathcal E,
 \qquad\overline q_i=e_i.                                  \tag{AP.11}
\]
Their first line and first two lines span precisely the
given ordered flag, since that flag has the stated distinct
inertial characters.  Consequently
$q_3\rho(g)q_1=0$ for every $g\in G_{E_v}$.
Put $u=\sum_iq_ie_i$.  It reduces to the identity, lies in
$\mathcal E$, and is a unit of that algebra: its inverse
is a polynomial in $u$ with coefficients in $\mathcal O$,
by Cayley--Hamilton and its unit determinant.
Also $ue_iu^{-1}=q_i$.  Applying (AP.8) and conjugating by
$\Pi_f(u)$ now diagonalizes its restriction to $G_{E_v}$;
the only possible off-diagonal term is the image of
$q_3\rho(g)q_1$, which is zero.  Thus the cohomology class,
not just its inertia restriction, vanishes at every $v\mid p$.
This argument does not require the global marked projectors
to be ordinary projectors or to commute with the local flag.

Away from $p$, $\rho$ is trivial on $E$-inertia, so
the class is unramified there.  The descended rational
class is zero locally at $p$ and unramified at every
$q\ne p$: restriction to the relevant $E$-decomposition
or inertia subgroup has kernel of order at most two in
the inflation--restriction calculation.  Positive-degree
cohomology of this quotient vanishes because $p>2$.
This also treats the ramified primes $q\mid D$.
We have proved (AP.7).

For the stated source application, \texttt{IsPtOf.lean}
gives the entrywise split residual representation, the
normalization and integrality.  Finite coefficients were
explained above.  \texttt{FOneCharU.lean:47--54} supplies
the Type~I admissibility predicate; \texttt{AfStmt.lean:29--32}
includes its level condition.  In the original
\texttt{BlochKatoSelmerVanishes.lean},
\texttt{IsFlagOn:37--39}, \texttt{LevelAt:102--107} and
\texttt{Level.above:111--114} give exactly the ordered
decomposition-group flag used here.  The finite characters
are trivial because \texttt{KillsMu 2}, defined at lines
66--67, is trivial when $\mu_2\subset\mathbf Q$.
The away-prime clause follows either from
\texttt{LevelAway:89--91} or directly from
\texttt{PtInertiaAway.lean:54--81}.  No new automorphy or
ordinary deformation theorem is needed for this application.
\end{proof}

\begin{corollary}[KT-ACTUAL-LOWER-LONG-FACTORIZATION]
Under the preceding theorem, if
$B_{1,\eta\omega^{2e}}\in\mathbf Z_p^\times$, then
\[
              \mathfrak a_{31}
                =\mathfrak a_{32}\mathfrak a_{21}
                =\mathfrak a_{21}^{2}.                     \tag{AP.12}
\]
If $\rho$ is absolutely irreducible, these are nonzero ideals,
so $v_\varpi(\mathfrak a_{31})=2v_\varpi(\mathfrak a_{21})$.
\end{corollary}
\begin{proof}
KT-AUXILIARY-UNIT-SELMER-BRIDGE, with $a=2e$, makes
the right side of (AP.7) zero, including after the finite
residue-field extension to $k$.  Hence the finite module
$C=\mathfrak a_{31}/\mathfrak a_{32}\mathfrak a_{21}$
has $\operatorname{Hom}_{\mathcal O}(C,k)=0$.
Thus $C/\varpi C=0$, and Nakayama gives $C=0$.
The second equality is (AP.3).  Absolute irreducibility
makes $F[\rho(H)]=M_3(F)$ by Burnside, so every block
ideal is nonzero.  The valuation identity follows.
\end{proof}

\paragraph{What the published specialization theorem does and does not add.}
Bella\"iche--Chenevier, \emph{The sign of Galois representations
attached to automorphic forms for unitary groups},
Compositio Math.\ 147 (2011), Proposition~2.4, pp.~1343--1344,
assumes a representation over a henselian DVR and
multiplicity-free absolutely irreducible constituents on
both fibers.  Its forward implication carries positive
signs of polarized generic constituents to the polarized
residual constituents.  Its converse additionally requires
the same number of irreducible constituents in both fibers.
The complete proof explicitly explains why an arbitrary
characteristic-zero polarization matrix need not be integral
and invertible on the given lattice.  Its argument uses
involution-stable idempotents, not unqualified reduction of
that matrix.

Our residual triple has three constituents, whereas an
irreducible rank-three generic point has one, and an outer
irreducible rank-two plus middle rank-one point has two.
That converse therefore does not infer the needed generic
sign from this residual triple.  The forward implication
could apply to a separately constructed good generic family,
but does not construct that family on every universal
Cayley--Hamilton branch.  Theorem~1.2 of the same paper has
an actual regular algebraic polarized cuspidal automorphic
input; it is not a statement about every pseudodeformation.
Here (AP.2) supplies the required sign and lattice directly.
Similarly, the hypotheses preceding Proposition~1.8.10 in
\emph{Families of Galois representations and Selmer groups},
\S1.8.7, pp.~56--58, require generic irreducibility in
every fraction-field factor.  That proposition cannot be
applied to the disputed $M_2\times F$ case merely because
the total dimension is three.

\paragraph{Remaining comparison boundary and prior work.}
The fixed $J$ identity was already used in the host, projector
and sparse-jet arguments.  The new step here is its direct
transport to the actual matrix-image extension module,
together with the ordinary local splitting proof and (AP.12).
This is distinct from KT-SELFDUAL-SOURCE-ODD-DESCENT:
conjugate self-duality does not assert $\rho^\vee\simeq\rho$.
It also does not supply the depth or sparse-support hypotheses
of KT-EQUIVARIANT-POINT-NOT-CUBIC-JET, or the adjacent
dimension bound in KT-RESIDUAL-CYCLICITY-SCOPE.

In particular, (AP.12) allows nonzero lower long entries
produced by products of adjacent entries.  It does not make
$\mathfrak a_{21}$ zero.  This is compatible with the earlier
mixed-path and higher-obstruction calculations.  It repairs
the long-root sign for the actual image and yields a precise
arithmetic factorization there, but excludes no remaining
primitive nonessentially-self-dual Type~I source.

A surjection from a universal Cayley--Hamilton algebra to
$\mathcal E$ gives no reverse injection of its long-root
module into the one in (AP.7).  Its kernel can carry
extension data.  To use the same repair for that universal
module one must separately prove a faithful compatible
matrix realization with a common symmetric pairing, or
prove that its kernel is invisible to the particular
finite-coefficient Hom module.  Neither premise is supplied
by (AP.2), by semisimple $R=T$, or by the inspected source
point construction.  The preprint's weight cone, ordinary
orders, numerical Selmer bounds and missing extra adjacent
congruence factor also remain as previously recorded.
Accordingly this is a repaired actual-image comparison,
classified as NONMOVEMENT for the main vanishing gate.

\paragraph{Arithmetic--geometric location.}
At each finite coefficient level, the arithmetic lifts in (AP.7)
define locally constant sheaves on
$\operatorname{Spec}\mathcal O_E[1/p]$.
The common perfect matrix $J$ pairs the integral sheaf with
its conjugate dual; the marked idempotents cut out the actual
matrix-image extension modules.  The target of (AP.7) has the
quadratic-twisted negative Tate coefficient
$k(\eta\omega^{-2e})$.  It is not the untwisted rational
coefficient in (FW.2).  The auxiliary unit annihilates this
particular strict extension group and yields (AP.12).
The remaining adjacent rational extension group is still the
strict Tate-sheaf group of (FW.2), whose dual is the integral
cohomology and $K$-group in (FW.1).  No rationalization of
these sheaves, no comparison of their generic traces, and
no additional assertion about an integral category of motives
is used to discard that remaining finite obstruction.

\paragraph{Exact source and review scope.}
At the pinned source commit, the underlying identities are
\texttt{UnitaryThreeReflectionRight.lean:44--85,283--289},
\texttt{ConjugateStarMatrix.lean:5--56}, and
\texttt{GradedModelDerivationConditionOfPtHyp.lean:93--107}.
The latter proves that the marked $\zeta$ is a unit.
\texttt{BaseRepGL3NormalisedOfPsiEqOne.lean:26--35} gives
the literal Type~II corner, and
\texttt{RhoKApply.lean:76--83} records the point predicates.
The matrix extension construction and its entire proof were
checked against Bella\"iche--Chenevier, Ast\'erisque 324
(2009), Theorem~1.5.5, pp.~38--39; their full
\S1.8.7 argument, pp.~56--58, and the entire published
2011 Proposition~2.4 proof, pp.~1343--1344, were read.
The 2011 main sign theorem is used only to delimit its
statement's scope here; its long automorphic proof is not
claimed freshly audited.  Berger--Klosin's full
Theorem~6.4/Proposition~6.7/Lemma~6.8 argument, pp.~32--36,
was reread for comparison.  The auxiliary-unit theorem
is an already checked input.  No Lean compilation or
new finite-degree $K$-group computation is claimed.
